
\documentclass[12pt]{article}

% Layout.
\usepackage[top=1in, bottom=0.75in, left=1in, right=1in, headheight=1in, headsep=6pt]{geometry}

% Fonts.
\usepackage{mathptmx}
\usepackage[scaled=0.86]{helvet}
\renewcommand{\emph}[1]{\textsf{\textbf{#1}}}

% TiKZ.
\usepackage{tikz, pgfplots}
\usetikzlibrary{calc,trees,positioning,arrows,fit,shapes,calc}
\pgfplotsset{compat = newest}
\tikzset{
  jumpdot/.style={mark=*,solid},
  excl/.append style={jumpdot,fill=white},
  incl/.append style={jumpdot,fill=black},
}
 
\pgfplotsset{my style/.append style={axis x line=middle, axis y line=
middle, xlabel={$x$}, ylabel={$y$}, axis equal }}

% Misc packages.
\usepackage{amsmath,amssymb,latexsym}
\usepackage{graphicx}
\usepackage{array}
\usepackage{xcolor}
\usepackage{multicol}

% Commands to set various header/footer components.
\makeatletter
\def\doctitle#1{\gdef\@doctitle{#1}}
\doctitle{Use {\tt\textbackslash doctitle\{MY LABEL\}}.}
\def\docdate#1{\gdef\@docdate{#1}}
\docdate{Use {\tt\textbackslash docdate\{MY DATE\}}.}
\def\doccourse#1{\gdef\@doccourse{#1}}
\let\@doccourse\@empty
\def\docscoring#1{\gdef\@docscoring{#1}}
\let\@docscoring\@empty
\def\docversion#1{\gdef\@docversion{#1}}
\let\@docversion\@empty
\makeatother

% Headers and footers layout.
\makeatletter
\usepackage{fancyhdr}
\pagestyle{fancy}
\fancyhf{} % Clears all headers/footers.
\lhead{\baselineskip 30pt
%\emph{\@doctitle\hfill\@docdate}
\emph{\@docdate\hfill\@doctitle}
\ifnum \value{page} > 1\relax\else\\
\emph{Name: \rule{3.5in}{1pt}\hfill \@docscoring}\fi}
\rfoot{\emph{\@docversion}}
\lfoot{\emph{\@doccourse}}
\cfoot{\emph{\thepage}}
\renewcommand{\headrulewidth}{0pt}%
\makeatother

% Paragraph spacing
\parindent 0pt
\parskip 6pt plus 1pt

% A problem is a section-like command. Use \problem{5} to
% start a problem worth 5 points.
\newcounter{probcount}
\newcounter{subprobcount}
\setcounter{probcount}{0}
\newcommand{\problem}[1]{%
\par
\addvspace{4pt}%
\setcounter{subprobcount}{0}%
\stepcounter{probcount}%
%\makebox[0pt][r]{\emph{\arabic{probcount}.}\hskip1ex}\emph{[#1 points]}\hskip1ex}
\makebox[0pt][r]{\emph{\arabic{probcount}.}\hskip1ex}\emph{[#1 \ifnum1=0#1\relax point\else points\fi]}\hskip1ex}
%\textbf{(#1 \ifnum1=0#1\relax thing\else things\fi
\newcommand{\thesubproblem}{\emph{\alph{subprobcount}.}}

% Subproblems are an enumerate-like environment with a consistent
% numbering scheme. 
% Use \begin{subproblems}\item...\item...\end{subproblems}
\newenvironment{subproblems}{%
\begin{enumerate}%
\setcounter{enumi}{\value{subprobcount}}%
\renewcommand{\theenumi}{\emph{\alph{enumi}}}}%
{\setcounter{subprobcount}{\value{enumi}}\end{enumerate}}

% Blanks for answers in normal and math mode.
\newcommand{\blank}[1]{\rule{#1}{0.75pt}}
\newcommand{\mblank}[1]{\underline{\hspace{#1}}}
\def\emptybox(#1,#2){\framebox{\parbox[c][#2]{#1}{\rule{0pt}{0pt}}}}

% Misc.
\renewcommand{\d}{\displaystyle}
\newcommand{\ds}{\displaystyle}
\def\bc{\begin{center}}
\def\ec{\end{center}}
\def\be{\begin{enumerate}}
\def\ee{\end{enumerate}}

\newcommand{\ans}{\rule{2 in}{.5pt}}
\newcommand{\bigans}[1]{\rule{#1 in}{.5pt}}

\doctitle{Math F251X: Quiz 1}
\docdate{Fall 2026}
\doccourse{UAF Calculus I}
\docversion{v-1}
\docscoring{\blank{0.8in} / 25}
\begin{document}
\textbf{Please circle your section:} \hfill 901 (Gossell)  \hfill 902 (Gossell) \hfill 903 (Graves-McCleary)\\

There are 6 questions worth 25 points on this quiz. No aids (book,
calculator, etc.)
are permitted.  {\bf Show all work for full credit.}

%Problem
\problem{6} Given functions $f(x)=x^3-3x+1$ and $g(x)=2^x$, compute the following and \textbf{simplify} your answers:
\begin{subproblems}
\item $f(2)+g(2)$ \\
  \vfill
  \quad \hspace{4in} \underline{\hspace{2in}}
\item $f(0)-g(0)$ \\
  \vfill
  \quad \hspace{4in} \underline{\hspace{2in}}
\item $g(f(1))$ \\
  \vfill
  \quad \hspace{4in} \underline{\hspace{2in}}
\end{subproblems}

\problem{3} Write an equation of the line passing through the points $(1,2)$ and $(-1,5)$.\\

\vspace{1.5in}

\problem{3} Find all $x$-intercepts for the function $f(x)=2x^2-7x+5$.

\vspace{1.8in}

\newpage

\problem{4} In the right triangle below, suppose $a=5$ and $c=\sqrt{34}$.

\begin{multicols}{2}

\begin{tikzpicture}
  % Triangle
  \draw (0,0) node[anchor=north]{$A$}
    -- (4,0) node[anchor=north]{$C$}
    -- (4,4.5) node[anchor=south]{$B$}
    -- cycle;

  % Side labels
  \node at (2,-0.5){$b$};
  \node at (4.4,2.3){$a$};
  \node at (1.5,2.3){$c$};

  % Angle markers 
  \draw (0.65,0) arc[start angle=0,end angle=71,radius=0.45];
  \draw (3.59,4.05) arc[start angle=200,end angle=270,radius=0.45];

  % Right angle marker at C
  \draw (4,0) ++(-0.4,0) -- ++(0,0.4) -- ++(0.4,0);

\end{tikzpicture}

\vspace{.1in}

Calculate the following:\\

$\sin(A)=$\\

\vspace{.1in}

$\cos(A)=$\\

\vspace{.1in}

$\tan(A)=$\\

\end{multicols}

\problem{3} Evaluate the following:

\begin{multicols}{3}
\begin{subproblems}
\item $\displaystyle \sin\left(\frac{\pi}{4}\right)=$

\item $\displaystyle \tan\left(\frac{5\pi}{6}\right)=$

\item $\displaystyle \sec\left(\pi\right)=$

\end{subproblems}
\end{multicols}

\vspace{.2in}

\problem{6} State the domain and range of the following functions. Use interval notation.\\
\begin{subproblems}
\item $h(x) = \ln(x)$ \\
  \vfill
   domain:  \underline{\hspace{2in}} \hfill range:  \underline{\hspace{2in}}\\
\item $\displaystyle{g(x) = 3\cos(x)}+1$ \\
  \vfill
   domain:  \underline{\hspace{2in}} \hfill range:  \underline{\hspace{2in}}\\
\item $f(x) = -\sqrt{x+1}-5$ \\
  \vfill
   domain:  \underline{\hspace{2in}} \hfill range:  \underline{\hspace{2in}}\\
\end{subproblems}


\end{document}
%%% Local Variables:
%%% mode: LaTeX
%%% TeX-master: t
%%% End:
